\documentclass[a4paper,11pt,dvipdfmx]{jsarticle}

\usepackage[top=20mm, bottom=25mm, left=20mm, right=20mm]{geometry}
\usepackage{url}
\usepackage{graphicx}
\usepackage{amssymb}
\usepackage{amsmath}
\usepackage{booktabs}
\usepackage{float}
\usepackage{multirow}

\title{今週の進捗報告\\[0.5em]\large デマンド目標値を考慮したEMS最適化モデルの\\評価方法の見直し}
\author{神戸大学大学院システム情報学研究科 修士1年\\山崎博之}
\date{2026年7月11日}

\begin{document}

\maketitle

\section{背景と問題意識}
エネルギーマネジメントシステム（EMS）における蓄電池充放電計画では，電力量料金だけでなく，契約電力に基づく基本料金を適切に考慮する必要がある．特に高圧契約では，契約電力は原則として当月を含む過去12か月の最大需要電力に基づいて決定されるため，ある時点で大きな買電ピークが発生すると，その影響が以後の月の基本料金に残り続ける．

6月の報告では，この料金制度を考慮するために，各月の基本料金を決定するデマンド目標値 $\bar{S}_{m}$ を導入した最適化モデルを定式化した．一方で，計算結果の比較方法については，従来手法と提案手法で基本料金の評価ルールが完全には揃っていなかった．そのため，得られたコスト差が「デマンド目標値を考慮した最適化モデルそのものの効果」を表しているのか，あるいは評価式の違いによって生じているのかを明確に切り分ける必要がある．

\section{従来の比較方法における課題}
現在の進捗報告書の表1では，従来手法の運用結果と提案手法の運用結果を比較している．しかし，両者では基本料金の計算方法が異なっている．

\begin{itemize}
    \item \textbf{従来手法の運用結果}：最適化対象年の中の最大買電電力 $S_{\mathrm{MAX}}$ によって基本料金を評価する．
    \item \textbf{提案手法の運用結果}：各月について，当月を含む過去12か月参照のデマンド目標値 $\bar{S}_{m}$ によって基本料金を評価する．
\end{itemize}

この比較では，従来手法と提案手法で運用スケジュールの作り方が異なるだけでなく，最終的なコスト評価に用いる基本料金の計算ルールも異なっている．したがって，評価のモノサシが揃っておらず，提案手法が実際の料金制度の下でどれだけ有効であるかを純粋に比較できていない．

実際に，6月の報告ではデマンド目標値を置いた方法が従来の方法に比べて年間トータルコストを500円低減する結果となった．しかし，この差額は非常に小さく，提案モデルの価値が十分に現れていない可能性がある．これは，従来手法側を従来の基本料金計算式で評価し，提案手法側を新しい基本料金計算式で評価しているため，両者の違いが最適化モデルの違いとして解釈しにくくなっていることが原因であると考えられる．

\section{評価方法の見直し}
今後の比較では，運用スケジュールの作り方は従来手法と提案手法で分けたまま，評価用のコスト計算式を実際の料金制度に統一する．すなわち，どちらの運用結果に対しても，各月の基本料金を過去12か月参照のデマンド目標値 $\bar{S}_{m}$ に基づいて再計算する．

評価用コストは次式で整理する．

\begin{equation}
    C_{\mathrm{eval}}
    =
    \sum_{m=1}^{12}
    \left(
    w_{\mathrm{basic}}^{\mathrm{month}} \bar{S}_{m}
    +
    \sum_{k \in \mathcal{K}^{(2)}_{m}}
    P^{\mathrm{BY}} S^{\mathrm{BY}}_{k} \Delta t
    \right).
\end{equation}

ここで，$w_{\mathrm{basic}}^{\mathrm{month}}$ は月額基本料金単価，$P^{\mathrm{BY}}$ は電力量料金単価，$S^{\mathrm{BY}}_{k}$ は時刻 $k$ における買電電力，$\Delta t$ は時間幅である．また，$\bar{S}_{m}$ は月 $m$ の基本料金を決定する契約電力であり，当月を含む過去12か月間の買電電力の最大値として評価する．

このように評価式を統一することで，比較対象は「従来モデルで作った運用スケジュール」と「デマンド目標値を考慮した提案モデルで作った運用スケジュール」の違いに限定される．

\section{具体的な実験手順}
見直し後の比較実験では，表\ref{tab:experiment_plan}に示す2つのパターンを計算する．

\begin{table}[H]
    \centering
    \caption{評価方法を統一した比較実験の手順}
    \label{tab:experiment_plan}
    \begin{tabular}{p{25mm}p{55mm}p{65mm}}
        \toprule
        比較パターン & 最適化時の扱い                                                             & 評価時の扱い \\
        \midrule
        パターンA（従来手法による運用）
               & デマンド目標値 $\bar{S}_{m}$ を考慮せず，従来の最適化モデルのまま充放電スケジュールを算出する．
               & 得られた買電電力系列を，過去12か月参照の $\bar{S}_{m}$ に基づく新しい基本料金計算式に当てはめて再評価する．               \\
        \midrule
        パターンB（提案手法による運用）
               & デマンド目標値 $\bar{S}_{m}$ を制約および目的関数に組み込み，実際の基本料金制度を意識して充放電スケジュールを算出する．
               & パターンAと同じく，過去12か月参照の $\bar{S}_{m}$ に基づく新しい基本料金計算式で評価する．                       \\
        \bottomrule
    \end{tabular}
\end{table}

パターンAでは，最適化モデル自体は従来手法のままでよい．ただし，その運用結果を評価するときに，新しい基本料金の計算式を付け足す．これにより，デマンド目標値を考慮しない運用が，実際の料金制度の下でどれだけ不利になるかを確認できる．

パターンBでは，最適化段階から $\bar{S}_{m}$ を考慮する．過去12か月の参照範囲内で高いピークが残る月では，ピーク抑制の効果がすぐには基本料金に反映されない可能性がある一方で，高いピークが参照範囲から外れる月以降では，新たなピークを立てないことが基本料金の低減に直結する．このような時間的な影響を最適化モデルに組み込む点が，提案手法の特徴である．

\section{再評価結果}
上記の方針に基づき，従来手法で得られた買電電力系列を，過去12か月参照のデマンド目標値 $\bar{S}_{m}$ に基づく新しい基本料金計算式で再評価した．提案手法については，6月報告で用いたデマンド目標値を考慮したLPモデルの出力を同じ評価式で集計した．

年間集計結果を表\ref{tab:reassessment_results}に示す．

\begin{table}[H]
    \centering
    \caption{評価方法を統一した年間コスト比較}
    \label{tab:reassessment_results}
    \begin{tabular}{lrrr}
        \toprule
        評価項目                   & パターンA（従来手法を再評価） & パターンB（提案手法）     & 差分（A--B） \\
        \midrule
        年間電力量料金 [円]            & 11,354,515         & 11,354,053       & 461      \\
        年間基本料金 [円]             & 4,815,077          & 4,815,039        & 38       \\
        年間トータルコスト [円]          & 16,169,592         & \textbf{16,169,092} & 499      \\
        最適化対象年の最大買電電力 [kW]     & 166.831            & 166.826          & 0.005    \\
        最大デマンド目標値 [kW]          & 166.831            & 166.831          & 0.000    \\
        \bottomrule
    \end{tabular}
\end{table}

評価式を新しい料金制度に統一しても，今回の条件では提案手法による年間トータルコストの低減額は約499円にとどまった．内訳を見ると，電力量料金の差が約461円，基本料金の差が約38円であり，基本料金の削減効果はほとんど表れていない．また，最大デマンド目標値は両手法とも166.831\,kWでほぼ一致している．

\section{結果の解釈と今後の課題}
今回の再評価で大きな差が出なかった主な理由は，過去実績年として用いた買電電力系列が，すでに従来LPモデルによってピーク抑制された運用結果であるためだと考えられる．すなわち，過去12か月参照の中に入る前年側のピークが166.831\,kW程度に抑えられており，提案手法が最適化対象年でわずかにピークを下げても，各月のデマンド目標値 $\bar{S}_{m}$ はほとんど低下しない．

この結果から，評価基準を統一するだけでは，提案手法の価値を十分に示せないことが確認された．提案手法の効果を明確に評価するためには，過去実績年に「デマンド目標値を考慮しない無邪気な運用」または実測に近い買電系列を与え，その高いピークが参照期間から外れる過程で，提案手法が新たなピークを抑制できるかを検証する必要がある．

今後は，過去実績年の買電系列として，蓄電池を用いない場合の正味需要や，デマンドを意識しない運用結果を用いるケースを追加する．その上で，年間電力量料金，年間基本料金，年間トータルコスト，月別デマンド目標値 $\bar{S}_{m}$，月内最大買電電力を比較し，デマンド目標値を考慮した最適化モデルの有用性を改めて評価する．

\end{document}
